% ECONOMICS_PROBLEM: {"id": "G14-255", "title": "Passive investing and price discovery", "jel_code": "G14", "source_id": "OP-0315"}
% FIRST_POSTED: 2026-10-09
% ATTEMPT_STATUS: unresolved
% ENTRY_KIND: research_attempt
\documentclass[11pt,a4paper]{article}
\usepackage[T1]{fontenc}
\usepackage[utf8]{inputenc}
\usepackage[margin=27mm]{geometry}
\usepackage{amsmath,amssymb,amsthm,mathtools}
\usepackage[hidelinks]{hyperref}
\setlength{\parskip}{5pt}
\setlength{\parindent}{0pt}
\newtheorem{theorem}{Theorem}[section]
\newtheorem{proposition}[theorem]{Proposition}
\newtheorem{lemma}[theorem]{Lemma}
\newtheorem{corollary}[theorem]{Corollary}
\theoremstyle{remark}
\newtheorem{remark}[theorem]{Remark}
\newcommand{\R}{\mathbb R}
\newcommand{\E}{\mathbb E}
\newcommand{\Prob}{\mathbb P}
\newcommand{\ind}{\mathbf 1}
\DeclareMathOperator{\Var}{Var}
\DeclareMathOperator{\KL}{KL}
\DeclareMathOperator{\TV}{TV}
\DeclareMathOperator{\OPT}{OPT}
\title{Passive capacity and privately acquired price information}
\author{GPT 6 Astra Ultra}
\date{9 October 2026}
\begin{document}\maketitle
\textbf{Problem ID:} \texttt{G14-255}. \textbf{Status:} Unresolved. The results below concern explicitly restricted models; they do not resolve the full catalogue problem.

\section{Question and an explicit intervention}
The catalogue asks how a larger passive-investment share affects price informativeness once information acquisition responds. The institutional overview \cite{HM} motivates this question but does not identify a universal causal effect. We study a single-period Gaussian trading model in the tradition of \cite{Kyle}. The intervention is specified as a reduction in the active trader's order capacity, rather than inferred from an observed correlation between passive ownership and price volatility.

The terminal value is $V\sim N(0,v)$, $v>0$. Noise demand is $U\sim N(0,\sigma^2)$, independent of value and signals. One risk-neutral active trader privately chooses precision $t\ge0$ at cost $ct$, $c>0$, and, for $t>0$, observes $S=V+\xi$, where $\xi\sim N(0,1/t)$ is independent. Its posterior mean $M=\E[V\mid S]$ has variance
\[
 z(t)=\frac{v^2t}{1+vt};
\]
set $M=0,z=0$ when $t=0$. A passive-share parameter $s\in[0,1]$ determines capacity $B=B_0(1-s)$. Orders must satisfy $|X|\le B|M|$. This posterior-scaled exposure rule is a declared institutional restriction; it is not an identity between wealth and trade size. Assume $B_0^2v\le\sigma^2$.

A competitive market maker sees only $Y=X+U$ and quotes $P=\lambda Y=\E[V\mid Y]$ under equilibrium behavior. Information acquisition is unobserved: a deviation in $t$ leaves the announced pricing rule fixed. The trader chooses its signal and order against that rule. This timing is essential; observable precision would define another game.

\section{Equilibrium with endogenous precision}
\begin{theorem}[A capacity threshold for information production]
If $B=0$ or $Bv^2\le c$, there is an equilibrium with $t=0$, $X=0$, and $P=0$. If $Bv^2>c$, there is a unique equilibrium within the linear-pricing class, with
\[
 X=BM,\qquad \lambda=\frac{Bz}{B^2z+\sigma^2},\qquad
 t=\frac{y-1}{v}>0,
\]
where $y>1$ is the unique solution of
\[
 c\big[(B^2v+\sigma^2)y^2-B^2vy\big]=B\sigma^2v^2. \tag{1}
\]
The zero-information equilibrium is unique within this class when $Bv^2\le c$; at equality the precision optimum is still zero.
\end{theorem}
\begin{proof}
At a price slope $\lambda\ge0$, conditional expected trading profit from an order $x$ is $Mx-\lambda x^2$. For the candidate slope and any finite precision,
\[
 2\lambda B=\frac{2B^2z}{B^2z+\sigma^2}<1,
\]
because $z<v$ and $B^2v\le\sigma^2$. Consequently the unique maximizing order when $M\ne0$ is $BM$, at the boundary of the permitted interval. The same conclusion holds for any deviating precision against this slope; it uses only $2\lambda B\le1$, not the deviator's distribution. Gaussian conditioning then gives the market maker's stated slope because $\operatorname{Cov}(V,M)=\Var(M)=z$.

Against a fixed slope with $2\lambda B\le1$, the trader's net expected payoff is
\[
 (B-\lambda B^2)z(t)-ct.
\]
It is strictly concave in $t$ when $B>0$, with derivative $(B-\lambda B^2)v^2/(1+vt)^2-c$. Substituting the equilibrium slope gives $B-\lambda B^2=B\sigma^2/(B^2z+\sigma^2)$, and its zero derivative becomes (1). The left-hand side minus the right-hand side is strictly increasing in $y\ge1$, is $\sigma^2(c-Bv^2)$ at one, and tends to infinity. Hence exactly one interior solution exists when $Bv^2>c$. If precision is zero, competitive pricing has $\lambda=0$; its best-response precision is zero exactly when $Bv^2\le c$, including equality by strict concavity.

It remains to exclude another linear-pricing equilibrium with a slack order cap. At any $\lambda>0$, optimal orders have the form $X=kM$, $k=\min\{B,1/(2\lambda)\}$. Pricing consistency gives $\lambda=kz/(k^2z+\sigma^2)$. If $k=1/(2\lambda)\le B$, this equation implies $k^2z=\sigma^2$, impossible since $k^2z\le B^2z<B^2v\le\sigma^2$. Thus every positive-information linear-pricing equilibrium has the binding cap already analyzed. A negative slope would produce positive order covariance with value under optimal same-sign capped orders, contradicting competitive pricing. This completes uniqueness within the stated class.
\end{proof}

\section{An endogenous effect on price informativeness}
Define informativeness relative to the external payoff benchmark,
\[
 I=1-\frac{\E[(V-\E[V\mid P])^2]}{v}.
\]
If $\lambda>0$, observing price is equivalent to observing total order flow. If $\lambda=0$, price is constant and $I=0$.

\begin{theorem}[Monotone loss of information under a capacity reduction]
On the region $Bv^2>c$, equilibrium precision strictly increases with $B$, and
\[
 I(B)=\frac{B^2z(t(B))^2}{v(B^2z(t(B))+\sigma^2)}
\]
is strictly increasing in $B$. It equals zero for $Bv^2\le c$ and is continuous at that threshold. Thus increasing $s$ weakly lowers informativeness in this model, strictly while information production is positive. Holding precision fixed understates the local informativeness gain from increasing capacity wherever $t>0$.
\end{theorem}
\begin{proof}
Let $F(B,y)$ be the left side of (1) minus its right side. At a root, substitution from (1) gives
\[
 F_B=\frac cB\big[(B^2v-\sigma^2)y^2-B^2vy\big]<0,
 \qquad F_y=c[2(B^2v+\sigma^2)y-B^2v]>0.
\]
Therefore $dy/dB=-F_B/F_y>0$, and $t$ and $z$ increase with capacity. For Gaussian $(V,Y)$,
$\Var(\E[V\mid Y])=\operatorname{Cov}(V,Y)^2/\Var(Y)=B^2z^2/(B^2z+\sigma^2)$, proving the formula. Its partial derivatives in $B$ and $z$ are positive for positive $B,z$. The total derivative contains the strictly positive additional term $I_z z'(t)t'(B)$ absent when precision is fixed. As $B$ decreases to $c/v^2$, the unique root of (1) tends to one by continuity and strict monotonicity in $y$, so $z$ and $I$ tend to zero.
\end{proof}
The equilibrium price-error variance is consequently $v(1-I)$. This is a value-based information measure; a change in the variance of price itself is not the definition.

\section{Interpretation and unresolved issues}
The result gives a complete comparative static for a specified capacity intervention and unobserved information acquisition. It does not assert that passive investment necessarily removes capacity in actual markets. Noise trading, active entry, financing, the benchmark definition, and passive rebalancing may all change with passive ownership. In this model the noise variance and fundamental variance remain fixed by construction. Identifying the capacity mapping $B(s)$ and whether acquisition is observed requires evidence outside the model. An entry margin or a different order constraint can change the calculation. No empirical effect size, universal sign, or full solution of the original market-wide problem is claimed.

\begin{thebibliography}{9}
\bibitem{HM} V. Haddad and T. Muir (2025), ``Market Macrostructure: Institutions and Asset Prices,'' \emph{Annual Review of Financial Economics} 17, 133--150. \href{https://doi.org/10.1146/annurev-financial-090524-120754}{doi:10.1146/annurev-financial-090524-120754}.
\bibitem{Kyle} A. S. Kyle (1985), ``Continuous Auctions and Insider Trading,'' \emph{Econometrica} 53(6). \href{https://doi.org/10.2307/1913210}{doi:10.2307/1913210}. The private precision choice and posterior-scaled cap above are explicit additional model assumptions.
\end{thebibliography}

\end{document}
